% GENERATED FILE. Edit canonical lessons, then run npm run generate:books.
% canonical-source-hash: fnv1a64:62844c8cc3967ae3
% canonical-lessons: KA-C227-garagasa, KA-C227-eni, KA-C227-antu, KA-C227-tarc, KA-C227-yantra

\chapter{Saw, Ladder, Glue, Torch, Machine}
\label{ch:ka-saw-ladder-glue-torch-machine}
\hlchaptermodality{227}

\begin{chapteropening}
\textbf{By the end of this chapter:} \emph{I can say a saw, a ladder, glue, a torch and a machine.}

Complete the last lesson of chapter 227: I can say a saw, a ladder, glue, a torch and a machine.
\end{chapteropening}

\section[\texorpdfstring{\kn{ಗರಗಸ} (garagasa)}{ಗರಗಸ (garagasa)}]{\kn{ಗರಗಸ} (garagasa) — a saw}

\label{lesson:KA-C227-garagasa}



\begin{quote}
Before the new one: say the Kannada for son-in-law, then the Kannada for sister-in-law.
\end{quote}

\subsection*{You'll want to know: \kn{ಗರಗಸ}}
\textbf{\kn{ಗರಗಸ}} — \emph{garagasa} — \textquotedblleft{}a saw\textquotedblright{}.

\begin{grammarlens}[title={What you've built}]
One more word for this chapter.
\end{grammarlens}

\subsection*{Your turn}
\begin{itemize}
\raggedright
  \item \emph{Say it:} \emph{garagasa}
  \item \emph{Say it:} \emph{garagasa}, once more
  \item \emph{Say it:} say \emph{garagasa}
  \item \emph{Recall:} say the Kannada for son-in-law, then the Kannada for sister-in-law, then say \emph{garagasa} again
\end{itemize}

\begin{tcolorbox}[breakable,colback=teal!4,colframe=teal!35!black,title={Before you move on}]
What does \kn{ಗರಗಸ} mean? (\textquotedblleft{}A saw\textquotedblright{}.) Say it once more.
\end{tcolorbox}
\section[\texorpdfstring{\kn{ಏಣಿ} (ēṇi)}{ಏಣಿ (ēṇi)}]{\kn{ಏಣಿ} (ēṇi) — a ladder}

\label{lesson:KA-C227-eni}



\begin{quote}
Before the new one: say the Kannada for sister-in-law, then the Kannada for a saw.
\end{quote}

\subsection*{You'll want to know: \kn{ಏಣಿ}}
\textbf{\kn{ಏಣಿ}} — \emph{ēṇi} — \textquotedblleft{}a ladder\textquotedblright{}.

\begin{grammarlens}[title={What you've built}]
One more word for this chapter.
\end{grammarlens}

\subsection*{Your turn}
\begin{itemize}
\raggedright
  \item \emph{Say it:} \emph{ēṇi}
  \item \emph{Say it:} \emph{ēṇi}, once more
  \item \emph{Say it:} say \emph{ēṇi}
  \item \emph{Recall:} say the Kannada for sister-in-law, then the Kannada for a saw, then say \emph{ēṇi} again
\end{itemize}

\begin{tcolorbox}[breakable,colback=teal!4,colframe=teal!35!black,title={Before you move on}]
What does \kn{ಏಣಿ} mean? (\textquotedblleft{}A ladder\textquotedblright{}.) Say it once more.
\end{tcolorbox}
\section[\texorpdfstring{\kn{ಅಂಟು} (aṇṭu)}{ಅಂಟು (aṇṭu)}]{\kn{ಅಂಟು} (aṇṭu) — glue}

\label{lesson:KA-C227-antu}



\begin{quote}
Before the new one: say the Kannada for a saw, then the Kannada for a ladder.
\end{quote}

\subsection*{You'll want to know: \kn{ಅಂಟು}}
\textbf{\kn{ಅಂಟು}} — \emph{aṇṭu} — \textquotedblleft{}glue\textquotedblright{}.

\begin{grammarlens}[title={What you've built}]
One more word for this chapter.
\end{grammarlens}

\subsection*{Your turn}
\begin{itemize}
\raggedright
  \item \emph{Say it:} \emph{aṇṭu}
  \item \emph{Say it:} \emph{aṇṭu}, once more
  \item \emph{Say it:} say \emph{aṇṭu}
  \item \emph{Recall:} say the Kannada for a saw, then the Kannada for a ladder, then say \emph{aṇṭu} again
\end{itemize}

\begin{tcolorbox}[breakable,colback=teal!4,colframe=teal!35!black,title={Before you move on}]
What does \kn{ಅಂಟು} mean? (\textquotedblleft{}Glue\textquotedblright{}.) Say it once more.
\end{tcolorbox}
\section[\texorpdfstring{\kn{ಟಾರ್ಚ್} (ṭārc)}{ಟಾರ್ಚ್ (ṭārc)}]{\kn{ಟಾರ್ಚ್} (ṭārc) — a torch}

\label{lesson:KA-C227-tarc}



\begin{quote}
Before the new one: say the Kannada for a ladder, then the Kannada for glue.
\end{quote}

\subsection*{You'll want to know: \kn{ಟಾರ್ಚ್}}
\textbf{\kn{ಟಾರ್ಚ್}} — \emph{ṭārc} — \textquotedblleft{}a torch\textquotedblright{}.

\begin{grammarlens}[title={What you've built}]
One more word for this chapter.
\end{grammarlens}

\subsection*{Your turn}
\begin{itemize}
\raggedright
  \item \emph{Say it:} \emph{ṭārc}
  \item \emph{Say it:} \emph{ṭārc}, once more
  \item \emph{Say it:} say \emph{ṭārc}
  \item \emph{Recall:} say the Kannada for a ladder, then the Kannada for glue, then say \emph{ṭārc} again
\end{itemize}

\begin{tcolorbox}[breakable,colback=teal!4,colframe=teal!35!black,title={Before you move on}]
What does \kn{ಟಾರ್ಚ್} mean? (\textquotedblleft{}A torch\textquotedblright{}.) Say it once more.
\end{tcolorbox}
\section[\texorpdfstring{\kn{ಯಂತ್ರ} (yantra)}{ಯಂತ್ರ (yantra)}]{\kn{ಯಂತ್ರ} (yantra) — a machine}

\label{lesson:KA-C227-yantra}



\begin{quote}
Before the new one: say the Kannada for glue, then the Kannada for a torch.
\end{quote}

\subsection*{You'll want to know: \kn{ಯಂತ್ರ}}
\textbf{\kn{ಯಂತ್ರ}} — \emph{yantra} — \textquotedblleft{}a machine\textquotedblright{}.

\begin{grammarlens}[title={What you've built}]
One more word for this chapter.
\end{grammarlens}

\subsection*{Your turn}
\begin{itemize}
\raggedright
  \item \emph{Say it:} \emph{yantra}
  \item \emph{Say it:} \emph{yantra}, once more
  \item \emph{Say it:} say \emph{yantra}
  \item \emph{Recall:} say the Kannada for glue, then the Kannada for a torch, then say \emph{yantra} again
\end{itemize}

\begin{tcolorbox}[breakable,colback=teal!4,colframe=teal!35!black,title={Before you move on}]
What does \kn{ಯಂತ್ರ} mean? (\textquotedblleft{}A machine\textquotedblright{}.) Say it once more.
\end{tcolorbox}
